Let $A\subseteq\mathbb{N}$ be a \emph{finite} nonempty set, and define
\[
B := \{ n\ge 1 : \exists a\in A\ \text{such that } a\mid n\},
\]
the set of positive integers divisible by at least one element of $A$.
For $x\ge 1$ define the counting function
\[F(x):=|B\cap[1,x]|\quad\text{and the density-initial-segment }\quad d(x):=\frac{F(x)}{x}.
\]
The problem asks whether the following inequality always holds:
\[
\forall m>n\ge \max(A):\qquad \frac{F(m)}{m} < 2\frac{F(n)}{n}.
\]
It also notes that the constant $2$ would be best possible, since for $A=\{a\}$, $n=2a-1$, $m=2a$ makes the ratio approach $2$ from below.
